\section{Methodology}

This part sits in the middle of the paper. It explains how each result was obtained and lays out the survey that
ranges across every dimension. The sections that follow give the selection argument itself. Here we describe the
shared machinery once, so that the later claims can be read without stopping to ask what the words
\emph{measured} or \emph{derived} stand for. The per-sector detail belongs to the longer methodology notes, and
the per-dimension data belongs to the appendices. What is gathered here is the common ground that every
computation in the paper rests on.

\subsection{The construction}

Everything begins with a single object, the cell-adjacency graph of a honeycomb. A \emph{cell-graph builder}
takes a Schl\"afli symbol and grows the honeycomb outward. Each cell is tracked as a group element, that is, as a
matrix in the symmetry group, and neighbours are reached by applying the face reflections. The traversal is a
breadth-first search whose cost scales with the facet count times the number of cells visited. This works
cleanly for the curved honeycombs that the paper cares about most. The flat references need a different route.
A separate \emph{Euclidean lattice builder} constructs the flat cases, the $D_4$ root lattice and the integer
lattices $\mathbb{Z}^n$, directly. The reason is geometric rather than a matter of taste. A flat honeycomb has
its cell centre at an ideal point, which is to say at infinity, and the hyperbolic builder stalls there. Building
the flat lattices by hand sidesteps that failure and gives clean controls.

\subsection{The probes}

A small set of probes turns the graph into numbers. Each one targets a different physical quantity, and together
they cover dimension, dynamics, gravity, and the internal symmetry content. We list them here and return to each
in the sectors where it does its work.

\begin{itemize}
  \item Breadth-first shells give the growth profile, which feeds dimension, curvature, cosmology, and the
  hierarchy of scales.
  \item The lazy random walk gives the spectral dimension through its return probability, which falls like
  $t^{-d/2}$.
  \item The Bethe-lattice cavity recursion gives the exact Green's function on the infinite tree, which feeds
  gravity and holography.
  \item The kernel polynomial method with a stochastic trace gives spectral quantities, such as the Skyrme sea
  energy, without ever diagonalising a matrix. It uses only matrix-vector products and scales onto a GPU.
  \item Representation theory reads off the directional coin's spinor and gauge content, which is where spin and
  the unified group come from.
  \item Variance-reduction methods, common random numbers and singular-part subtraction, cancel the
  finite-patch and boundary contamination that would otherwise corrupt the gravity law.
  \item WebGPU carries the large-scale dynamics, the $24$-direction lattice gas, the bulk wave, the skyrmion,
  and the kernel polynomial spectra.
\end{itemize}

\subsection{The hyperbolic-bulk toolkit}

Hyperbolic patches present a sharp practical problem. A finite ball in $H^n$ is almost all boundary, around $99$
percent of its cells touch the edge, so a naive bulk measurement is swamped by surface effects. A dedicated
toolkit, set out in full in the appendix, removes the confound by never trusting a raw finite patch. The Bethe
recursion works on the infinite tree directly, so there is no finite patch to corrupt. Closed manifolds remove
the boundary altogether. Hyperbolic band theory puts a momentum space on a closed manifold, which makes the
spectrum tractable. Tensor networks contract at polynomial cost. The kernel polynomial method reads spectra from
matrix-vector products alone. Bulk-local observables are designed to ignore the surface. The point of the whole
kit is that any quantity claimed for the bulk has been measured by a method that the boundary cannot reach.

\subsection{Reporting conventions}

Every result in the paper carries one of four tags. \emph{Derived} means it follows by argument from the base.
\emph{Measured} means it was computed on the substrate. \emph{Predicted} means the framework implies it but it has
not yet been checked. \emph{Open} means the question is genuinely unsettled. We hold ourselves to two further
commitments, and the appendix audits both.

\begin{principle}[Tag every claim]
No statement about the substrate is presented without one of the four tags. A measurement that was later found to
be confounded, such as the early gravity exponents taken on small patches, is flagged as confounded and
superseded by the result from the correct method, rather than quietly dropped.
\end{principle}

The second commitment is the nothing-added discipline already set out in Principle~\ref{prin:nothing-added}. No
computation in the paper introduces a new term, field, bias, or rule beyond the stated base in order to make a
result come out. When a phenomenon does not emerge from the base, the negative is reported. The base is fixed
before the experiments run, and it stays fixed.

These two rules are what let the selection argument carry weight. If the experiments could be tuned, a match
with known physics would prove nothing. Because the base is frozen and every claim is tagged, a match is
information and a failure is information, and the reader can tell which is which.
